\documentclass[10pt,a4paper]{article}

\usepackage[a4paper,top=12mm,bottom=12mm,left=13mm,right=13mm]{geometry}
\usepackage[T1]{fontenc}
\usepackage[utf8]{inputenc}
\usepackage{lmodern}
\usepackage{microtype}
\usepackage[dvipsnames]{xcolor}
\usepackage{enumitem}
\usepackage{tabularx}
\usepackage{array}
\usepackage[hidelinks]{hyperref}

\definecolor{ink}{HTML}{18212D}
\definecolor{muted}{HTML}{607080}
\definecolor{rule}{HTML}{C9D2DC}
\definecolor{accent}{HTML}{0B6F78}

\pagestyle{empty}
\setlength{\parindent}{0pt}
\setlength{\parskip}{0pt}
\renewcommand{\familydefault}{\sfdefault}
\renewcommand\labelitemi{\textcolor{accent}{\raisebox{0.25ex}{\scriptsize$\bullet$}}}
\setlist[itemize]{leftmargin=12pt,itemsep=2.35pt,topsep=2.8pt,parsep=0pt,partopsep=0pt}

\newcommand{\sectiontitle}[1]{%
  \vspace{9pt}
  {\small\bfseries\color{accent}\MakeUppercase{#1}}\par
  \vspace{2.7pt}{\color{rule}\hrule height 0.55pt}\vspace{5.2pt}
}
\newcommand{\sideitem}[2]{%
{\bfseries #1}\par{\small\color{muted}#2}\par\vspace{6pt}
}
\newcommand{\role}[3]{%
  {\bfseries #1}\hfill{\small\color{muted}#2}\par
  {\small\color{muted}#3}\par
}
\newcommand{\tightrole}[3]{%
  {\bfseries #1}\hfill{\small\color{muted}#2}\par
  {\small\color{muted}#3}\par\vspace{2pt}
}

\begin{document}
\fontsize{10.15}{12.25}\selectfont
\color{ink}

{\LARGE\bfseries Alexander Belik}\hfill{\small Dublin, Ireland}\\[-1pt]
{\large\color{accent} Theoretical Physics \,|\, Quantum Thermodynamics \,|\, Scientific Computing}\hfill{\small +353 85 256 4027}\\[3pt]
{\small
\href{mailto:abelik1@outlook.com}{abelik1@outlook.com}
\quad|\quad
\href{https://www.linkedin.com/in/alexander-belik/}{linkedin.com/in/alexander-belik}
\quad|\quad
\href{https://github.com/ABelik1}{github.com/ABelik1}
}
\vspace{7pt}

{\color{rule}\hrule height 0.8pt}
\vspace{8pt}

\begin{minipage}[t][0.835\textheight][t]{0.315\linewidth}
\raggedright

\sectiontitle{Profile}
Fourth-year Theoretical Physics student at Trinity College Dublin, focused on quantum thermodynamics, many-body dynamics, and research software.
\vspace{4pt}

I build reproducible Python and PyTorch tools that connect analytic models, simulation, machine learning, and experimental data workflows.

\sectiontitle{Education}
\sideitem{B.A. Theoretical Physics}{Trinity College Dublin, expected 2026}
{\small Quantum mechanics, statistical mechanics, quantum field theory, and computational physics.}\par\vspace{5pt}
\sideitem{Leaving Certificate}{The Dublin Academy of Education, 2020--2022}

\sectiontitle{Core Physics}
{\small
Quantum thermodynamics\par
Many-body dynamics\par
Gibbs-preserving maps\par
Density matrices and channels\par
Semidefinite programming\par
Spectroscopy workflows\par
Computational chemistry\par
}

\sectiontitle{Technical Toolkit}
{\small
Python, NumPy, SciPy\par
PyTorch, QuTiP, Optuna\par
Matplotlib, Mathematica\par
C/C++ simulations\par
Docker, SQLite\par
Git/GitHub, LaTeX\par
}

\sectiontitle{Research Practice}
{\small
Numerical simulation\par
Data analysis and diagnostics\par
Hyperparameter tuning\par
Reproducible pipelines\par
Technical writing\par
}

\vfill
\sectiontitle{Languages}
{\small English -- fluent\par Russian -- intermediate\par Spanish -- basic}

\end{minipage}
\hfill
\begin{minipage}[t][0.835\textheight][t]{0.645\linewidth}
\raggedright

\sectiontitle{Research Experience}

\role{Local Thermality from Global Athermality}{2025--2026}{Capstone thesis, supervisor: Felix Binder}
\begin{itemize}
  \item Studied locally thermal bipartite quantum states with Gibbs marginals, isolating correlations as the remaining thermodynamic resource.
  \item Characterized the locally thermal set as an affine slice of the positive semidefinite cone.
  \item Derived a nine-parameter two-qubit normal form and a Schur-complement positivity criterion.
  \item Analysed diagonal slices, locally thermal rays, commuting qutrit subclasses, and SDP convertibility tests using \(D(\rho\Vert\Gamma)=I(A:B)_\rho\) on the locally thermal manifold.
\end{itemize}
\vspace{2pt}

\role{Research Assistant, CRANN Institute}{Jun--Sep 2025}{Trinity College Dublin}
\begin{itemize}
  \item Designed Echo State Network reservoir models to learn quantum time-evolution data from simulated spin-chain observables.
  \item Built a PyTorch ESN toolkit with configurable reservoirs, mixed precision, ridge-regression solvers, Optuna tuning, and structured diagnostics.
  \item Paired the learning workflow with a Heisenberg-chain simulator and conservation checks for norm, energy, and magnetization.
\end{itemize}
\vspace{2pt}

\role{Experimental / Computational Research Assistant}{May--Sep 2024}{Trinity College Dublin}
\begin{itemize}
  \item Created Python experiment-control and data-acquisition tooling for spectroscopy measurements on ferrofluids and liquid crystals.
  \item Automated structured logging and repeatable measurement workflows; related work: \href{https://arxiv.org/abs/2504.19633}{arXiv:2504.19633}.
\end{itemize}
\vspace{2pt}

\role{DFT Hamilton / PySCF Work}{2024--2025}{Computational chemistry}
\begin{itemize}
  \item Built a Dockerized PySCF workflow for H2 RHF/DFT, open-shell Li UHF, molecular examples, spin-correlation tests, endpoint checks, and custom functional comparisons.
\end{itemize}

\sectiontitle{Teaching And Technical Work}

\tightrole{Teaching Assistant}{Oct 2025--Present}{Trinity College Dublin}
\begin{itemize}
  \item Lead weekly tutorials and problem-solving sessions aligned with lectures.
  \item Grade assignments and exams with clear, constructive feedback.
\end{itemize}
\vspace{2pt}

\tightrole{Software Engineer, EdTech}{2024--2025}{Grinds360}
\begin{itemize}
  \item Built Python automation tools for content and workflow pipelines.
  \item Contributed to product features and UX iteration with technical and non-technical stakeholders.
\end{itemize}

\vfill
\sectiontitle{Selected Research Software}
{\small
\textbf{EchoState and Heisenberg Chain:} physics-informed reservoir-computing toolkit for quantum dynamics prediction.\\[2pt]
\textbf{Spectroscopy Automation:} experiment-control and acquisition tooling for cleaner lab data capture.\\[2pt]
\textbf{PySCF Environment:} reproducible Docker workflow for computational chemistry experiments.
}

\end{minipage}

\end{document}
